\section{Data and methods}

\subsection{Sample and series}

The catalogue is the frozen version of Table~2 of Baliunas et al.\ \cite{baliunas1995} (\NCatalog{} stars), whose classes are Flat (S-index scatter below about 1--1.5\,\% of the mean), Cycle (a cycle period with a false-alarm class; we use the ``excellent'' ones), Var (variable without a periodic signal) and Long (variability on timescales longer than 25 years). Series come from three sources: the \NBaumStars{} stars of Baum et al.\ \cite{baum2022} (Table~1; points flagged as outliers are removed, points flagged both as one-day and outlier are kept), the \NVTwoStars{} stars of the HK Project v2001 files, and the calibrated solar S index of Egeland et al.\ \cite{egeland2017} (Table~4, 1966--2016, a composite of several instruments). Calibration stars, two binaries and stars with no series in any source were excluded before any indicator was computed. Our binarity rule leaves five clean Flat stars, of which four have a series. The stars analyzed are \NGroupFlat{} Flat, \NGroupCycle{} Cycle, \NGroupVar{} Var and \NGroupLong{} Long, plus the Sun, which is treated as exploratory.

For the temporal indicators each star is reduced to annual means of S, the mean of all observations in a calendar year. Only the Mount Wilson era is used, so no segment crosses the change to Keck; the change between the two Mount Wilson photometers inside that era is not corrected, and the solar composite combines five instruments as calibrated by Egeland et al. Gaps of one or two missing years are filled by linear interpolation and the filled points are marked; gaps of three or more years split the series into segments. Indicators are computed only in segments of at least \MinSeg{} years, so that a window of 50\,\% has at least 12 points and 13 windows. Segments of \MinSeg--\LowMax{} years are labeled low precision and their results exploratory; segments of \AdeqMin{} years or more are labeled adequate. \NStarsRun{} stars have such a segment (\LMin--\LMax{} years) and so does the Sun (\NSunYears{} years); \NNotRun{} further stars have a Mount Wilson segment shorter than \MinSeg{} years and receive a description only.

\subsection{Indicators and detrending}

In each segment we compute the lag-1 autocorrelation (AR(1) coefficient) and the variance in a sliding window, and the Kendall $\tau$ of each against window index (no tie correction; the same computation is used for the data and the surrogates). The window is a fraction $f$ of the segment length $L$, $w = \mathrm{round}(fL)$, with $f = 0.25$, 0.375, 0.5, 0.625 and 0.75; $f = 0.5$ is the default. The series is analyzed in four configurations: the original series, and the residual of a Gaussian smooth with $\sigma = 1$, 3 or 5 years (edges handled by repeating the end values). A Gaussian smooth of width $\sigma$ keeps a fraction $\exp[-\tfrac{1}{2}(2\pi\sigma/P)^2]$ of the amplitude of a cycle of period $P$; the residual is thus a high-pass version of the series, not the series without its cycle. A segment is flagged for an indicator when $\tau > 0$ with $p < 0.05$ in at least three of the four configurations. H2 requires this for both indicators.

\subsection{Significance}

The $p$ value of each $\tau$ comes from \NSurr{} surrogate series from an ARMA($p,q$) model with $p,q \le 2$ chosen by AIC, fitted by maximum likelihood to the series (or residual) being tested, with Gaussian innovations and a burn-in of 100 steps. Each surrogate receives the same gap positions and the same linear interpolation as the data. We report $p = P(\tau^{*} \ge \tau)$, where $\tau^{*}$ is the surrogate statistic, with the correction $(k+1)/(N+1)$ for $k$ exceedances in $N$ surrogates; this is the tail in the direction of critical slowing down. The size of the procedure was checked on synthetic AR(1) series with no trend (600 series of 30 years, three interpolated points, coefficient 0.5), with the acceptance range of \SizeLo--\SizeHi\,\% for a nominal 5\,\% fixed beforehand. Power was examined in an earlier simulation, before any stellar series was analyzed, in which the AR(1) coefficient either stayed at 0.5 or rose linearly from 0.3 to 0.7 (600 series per length, window 50\,\%, no cycle); the quantity reported is the area under the ROC curve (AUC) of $\tau$ for trend against null.

\subsection{State comparison for HD~166620}

For H1 we use the raw (not annual) series of Luhn et al.\ \cite{luhn2022}, in the three states they define: minimum~1 (1971--1975), minimum~2 (1987--1991) and the grand minimum (from 2004). The variability $\sigma_S$ of a state is the standard deviation of the interval with the instrumental noise variance removed in quadrature, the noise being \NoiseMW\,\% of the mean S for Mount Wilson and \NoiseKeck\,\% for Keck as quoted by Luhn et al.\ (method A, primary). As a sensitivity check, method B takes the mean over seasons of the same corrected value. Variances are compared with a two-sided $F$ test, with raw and with corrected values; the test treats observations as independent, which they are not (Sect.~\ref{sec:disc}).

\subsection{R$'_{\rm HK}$}

S is also converted to $R'_{\rm HK}$ with the recipe and color correction of Noyes et al.\ \cite{noyes1984}, using $B\!-\!V$ from the California Legacy Survey V catalogue \cite{isaacson2024} when the star is in it and from Baliunas et al.\ otherwise. The photospheric term of that calibration is valid for $0.44 < B\!-\!V < 0.82$; stars outside are marked in Table~\ref{tab:seg}. For a given $B\!-\!V$ the conversion is linear in S, so $\tau$ is unchanged and the run serves as a check.

\subsection{Group comparisons}

Groups are compared per star with the lag-1 autocorrelation and the normalized variance of the longest segment. The primary pair is Var and Long together against Flat, lag-1 autocorrelation of the $\sigma = 1$ residual, one-sided Mann--Whitney test, fixed before the tests were read. The remaining pairs (five per metric and configuration) are corrected with the Holm procedure within each metric and configuration block; Holm over all 30 tests is also reported (see below). The groups have four to eight stars, so the smallest attainable one-sided $p$ is \MinPCycleFlat{} (eight against four) or \MinPFourFour{} (four against four).

\subsection{Analysis log}

Fixed before the stellar run: the sample and groups, the segment rule, windows, detrending, surrogate design, H1 and H2, and the size check with its acceptance range. Decided after the run, without changing any computed value: the tail of $p$ (upper tail; the lower tail is also computed); keeping the points flagged both as one-day and outlier, which affects two stars (HD~190360 and HD~217014; with the opposite choice HD~190360 would fall below the minimum length); and the family of the Holm correction (blocks of five instead of all 30 tests). Both Holm families are reported.

\subsection{Software}

Python with NumPy, pandas, SciPy and statsmodels; the pipeline, the seeds (base \SeedBase) and the single official run are in the repository.
